\documentclass[10pt,a4paper]{article}
\usepackage[utf8]{inputenc}
\usepackage[T1]{fontenc}
\usepackage[french]{babel}
\usepackage{lmodern}
\usepackage[margin=2cm]{geometry}
\usepackage{graphicx}
\usepackage{booktabs}
\usepackage{amsmath,amssymb}
\usepackage{float}
\usepackage{caption}
\usepackage{microtype}
\usepackage{xcolor}
\usepackage[colorlinks=true,linkcolor=blue!50!black,urlcolor=blue!50!black,hypertexnames=false]{hyperref}
\usepackage{CJKutf8}

\newcommand{\zh}[1]{\begin{CJK}{UTF8}{gbsn}#1\end{CJK}}

\captionsetup{font=small,labelfont=bf}
\setlength{\parskip}{0.4em}

\begin{document}

%======================== PAGE DE GARDE ========================
\begin{titlepage}
\centering
\vspace*{1cm}
{\large Projet de Machine Learning\par}
\vspace{2cm}
\rule{\textwidth}{1.2pt}\\[0.6cm]
{\Huge\bfseries Classification d'images \\[0.3cm] du jeu de données Chinese MNIST\par}
\vspace{0.4cm}
{\Large Random Forest, réseau de neurones (MLP), SVM et CNN\par}
\rule{\textwidth}{1.2pt}
\vspace{2.5cm}

{\large\textbf{Auteurs}\\[0.3cm]
Léo \textsc{Halepidis}\\
Maxence \textsc{Lion}\par}
\vspace{1.5cm}

{\large\textbf{Professeur}\\[0.3cm]
Xavier \textsc{Siebert}\par}
\vspace{0.8cm}

{\large\textbf{Assistant}\\[0.3cm]
Timothy \textsc{Baeckelant}\par}
\vspace{1.5cm}
\end{titlepage}

%======================== SOMMAIRE ========================
\pagenumbering{gobble}
\tableofcontents
\clearpage
\pagenumbering{arabic}

%======================== INTRODUCTION ========================
\section{Introduction}

Le jeu de données \emph{Chinese MNIST} rassemble 15\,000 images en niveaux de gris de nombres manuscrits chinois, réparties en 15 classes équilibrées de 1\,000 images. On y retrouve les chiffres de \zh{零} (0) à \zh{九} (9), suivis des puissances de dix : \zh{十} (10), \zh{百} (100), \zh{千} (1\,000), \zh{万} (10\,000) et \zh{亿} ($10^8$).

Comparer plusieurs classifieurs vus au cours sur ces 15 classes constitue le fil conducteur de ce travail. Trois modèles de \texttt{scikit-learn} sont d'abord évalués sur les images aplaties en vecteurs 1D de 1\,024 pixels : un Random Forest, un réseau multicouche (\texttt{MLPClassifier}) et un SVM. Mais écraser une image en une seule ligne de pixels brise les relations de voisinage 2D. Un réseau convolutif (CNN, section~\ref{sec:bonus}) prenant directement la matrice $32\times32$ en entrée a donc été ajouté en bonus pour dépasser cette limite.

La section~\ref{sec:eda} présente d'abord le jeu de données et la section~\ref{sec:meth} détaille le protocole expérimental. Les résultats des quatre modèles sont ensuite regroupés à la section~\ref{sec:res}, suivis d'une analyse des confusions entre caractères proches (section~\ref{sec:crit}).

%======================== EDA ========================
\section{Exploration des données}\label{sec:eda}

Dans le fichier \texttt{chinese\_mnist.csv}, chaque ligne associe une image (\texttt{input\_<série>\_<échantillon>\_<code>.jpg}, de taille $64\times64$ pixels) à son code de classe (de 1 à 15), sa valeur numérique et son caractère chinois. Comme les échantillons proviennent de plusieurs scripteurs, l'épaisseur du trait, le centrage et l'inclinaison varient fortement d'une image à l'autre.

Chaque classe compte exactement 1\,000 exemples : le jeu est parfaitement équilibré. L'\emph{accuracy} globale et le $F_1$-score macro seront utilisés conjointement pour évaluer nos modèles.

\begin{table}[H]
\centering
\caption{Le jeu de données et notre découpage.}
\begin{tabular}{ll}
\toprule
Nombre d'images & 15\,000 (15 classes $\times$ 1\,000) \\
Taille d'origine & $64\times64$ pixels, niveaux de gris \\
Taille utilisée & $32\times32$ pixels (1\,024 variables en 1D, ou $32\times32\times1$ en 2D) \\
Normalisation & pixels divisés par 255 \\
Entraînement / test & 12\,000 / 3\,000 (stratifié, 200 images de test par classe) \\
\bottomrule
\end{tabular}
\end{table}
\begin{figure}[H]
    \centering
    \includegraphics[width=0.7\textwidth]{apercu_dataset.png}
    \caption{Visualisation d'un échantillon pour chacune des 15 classes du jeu de données Chinese MNIST.}
    \label{fig:apercu}
\end{figure}
Un simple coup d'œil aux images de la figure~\ref{fig:apercu} laisse déjà deviner certaines erreurs à venir : \zh{二} (2) et \zh{三} (3) ne se distinguent que par une barre horizontale, tandis que \zh{千} (1\,000) reprend la croix de \zh{十} (10) avec une courte barre oblique au sommet. De même, \zh{九} (9) et \zh{亿} ($10^8$) partagent une terminaison en crochet très proche (voir section~\ref{sec:crit}).

%======================== METHODOLOGIE ========================
\section{Méthodologie}\label{sec:meth}

\subsection{Prétraitement et séparation des données}
Chaque image est convertie en niveaux de gris et redimensionnée en $32\times32$ pixels par interpolation de Lanczos(pour préserver la netteté des traits), puis ses pixels sont normalisés entre 0 et 1 (division par 255). Travailler en $64\times64$ (4\,096 variables) rendait l'entraînement beaucoup trop lent ; passer en $32\times32$ divise la dimension par quatre sans détruire la forme générale des caractères. Pour les trois modèles de base, chaque image est aplatie en un vecteur $x = (x_1, x_2, \dots, x_d)^T \in \mathbb{R}^d$ avec $d = 1\,024$. Pour le CNN, ce vecteur est remodelé en tenseur 2D de taille $32\times32\times1$.

Le jeu de données est coupé en 80\,\% d'entraînement (12\,000 images, soit 800 par classe) et 20\,\% de test (3\,000 images, soit 200 par classe) avec l'option \texttt{stratify=y}. Les 3\,000 images de test sont isolées dès le début du script et ne servent qu'à l'évaluation finale.

\subsection{Reproductibilité (Graine aléatoire)}
La graine aléatoire est fixée à \texttt{seed = 42} dès le début du script (\texttt{scikit-learn}, \texttt{NumPy} et \texttt{Keras} via \texttt{keras.utils.set\_random\_seed}). Tous les algorithmes travaillent ainsi sur le même découpage. Comme nous n'avons testé qu'une seule graine, les scores affichés peuvent tout de même bouger de quelques dixièmes de point avec un autre tirage initial.

\subsection{Choix des hyperparamètres et stratégie de calcul}
Tester toutes les grilles sur les 12\,000 images prenait beaucoup trop de temps sur nos machines. C'était vrai pour le SVM — qui ajuste $\binom{15}{2} = 105$ classifieurs binaires à chaque pli — comme pour le CNN sur CPU. Un sous-ensemble stratifié de 3\,000 images d'entraînement (200 par classe) a donc été isolé pour régler les hyperparamètres sans jamais toucher au jeu de test (\emph{Data Leakage}).

Sur ces 3\,000 images, les trois modèles de \texttt{scikit-learn} passent par une validation croisée à 3 plis (\texttt{GridSearchCV}). Pour le CNN, quelques essais rapides ont d'abord écarté les valeurs trop lentes ($\eta=0{,}005$), avant de lancer une grille de 12 combinaisons sur 12 époques avec une découpe 80/20 (\texttt{X\_sub\_tr} / \texttt{X\_sub\_val}). Une fois la meilleure configuration retenue, chaque modèle est réentraîné sur les 12\,000 images complètes (tableau~\ref{tab:grilles}).

\begin{table}[H]
\centering
\caption{Grilles finales testées sur le sous-ensemble de 3\,000 images. La valeur retenue est en gras.}
\label{tab:grilles}
\small
\begin{tabular}{p{2.3cm}p{6.7cm}p{2.9cm}}
\toprule
\textbf{Modèle} & \textbf{Grille} & \textbf{Combinaisons} \\
\midrule
Random Forest & \texttt{n\_estimators} $\in\{100,\mathbf{300}\}$, \texttt{max\_depth} $\in\{\mathbf{None},20\}$, \texttt{max\_features} $\in\{\texttt{sqrt},\mathbf{0{,}1}\}$ & $2^3=8$ \\
\addlinespace
MLP & \texttt{hidden\_layer\_sizes} $\in\{(128),(256),\mathbf{(256,128)}\}$, \texttt{alpha} $\in\{\mathbf{10^{-4}},10^{-2}\}$, \texttt{learning\_rate\_init} $\in\{10^{-3},\mathbf{3\cdot10^{-3}}\}$ & $3\cdot2^2=12$ \\
\addlinespace
SVM & RBF : $C\in\{1,\mathbf{10},100\}$, $\gamma\in\{\mathbf{scale},10^{-3}\}$ ; linéaire : $C\in\{10^{-2},1\}$ & $6+2=8$ \\
\addlinespace
CNN (2D) & \texttt{lr} $\in\{0{,}01,\mathbf{0{,}02}\}$, \texttt{batch\_size} $\in\{\mathbf{32},64\}$, \texttt{dropout} $\in\{0{,}0,0{,}1,\mathbf{0{,}2}\}$ & $2\cdot2\cdot3=12$ \\
\bottomrule
\end{tabular}
\end{table}

Le tableau~\ref{tab:params} récapitule les paramètres retenus à l'issue de cette recherche et ceux fixés dans le code.

\begin{table}[H]
\centering
\caption{Hyperparamètres de chaque modèle : trouvés par le GridSearch ou fixés dans le code.}
\label{tab:params}
\small
\begin{tabular}{p{2.3cm}p{5.0cm}p{5.4cm}}
\toprule
\textbf{Modèle} & \textbf{Trouvés par le GridSearch} & \textbf{Fixés dans le code} \\
\midrule
Random Forest & \texttt{n\_estimators}=300, \texttt{max\_depth}=None, \texttt{max\_features}=0,1 & \texttt{random\_state}=42 \\
\addlinespace
MLP & \texttt{hidden\_layer\_sizes}=(256,128), \texttt{alpha}=0,0001, \texttt{learning\_rate\_init}=0,003 & \texttt{max\_iter}=200, \texttt{early\_stopping}=True, \texttt{random\_state}=42 \\
\addlinespace
SVM & noyau RBF, $C=10$, $\gamma$=scale & \texttt{random\_state}=42 \\
\addlinespace
CNN (2D) & \texttt{lr}=0,02, \texttt{batch\_size}=32, \texttt{dropout}=0,2 & 2 blocs \texttt{Conv2D} (32 et 64 filtres $3\times3$) + \texttt{MaxPooling2D} ($2\times2$), \texttt{Dense}(128), sortie \texttt{Softmax}(15), SGD (\texttt{momentum}=0,9), 25 époques \\
\bottomrule
\end{tabular}
\end{table}

\subsection{Les trois modèles de base (sur vecteurs 1D)}
Le \textbf{Random Forest} agrège par vote majoritaire 300 arbres de décision entraînés sans limite de profondeur sur des tirages bootstrap ; à chaque nœud, la coupure est tirée parmi 10\,\% des pixels (\texttt{max\_features=0.1}, soit une centaine de variables). Côté \textbf{MLP}, l'architecture retenue comporte deux couches cachées de 256 et 128 neurones avec une régularisation L2 ($\alpha=10^{-4}$) et l'option \texttt{early\_stopping=True}, qui réserve 10\,\% du train pour arrêter l'apprentissage dès que le score stagne. Quant au \textbf{SVM}, sa stratégie un-contre-un entraîne 105 classifieurs binaires : le \emph{GridSearch} a nettement préféré le noyau RBF ($C=10$, $\gamma=\texttt{scale}$) au noyau linéaire.

\subsection{Modèle bonus : Réseau de Neurones Convolutif (CNN)}\label{sec:bonus}
\paragraph{Rappel théorique du cours sur les réseaux de neurones.}
Dans un réseau multicouche de $L$ couches ($\ell = 1, \dots, L$) tel qu'étudié au cours, chaque neurone $j$ de la couche $(\ell)$ calcule une somme pondérée des sorties $a^{(\ell-1)}_k$ de la couche précédente avec les poids $w^{(\ell)}_{jk}$ et le biais $b^{(\ell)}_j$, puis applique une fonction d'activation $\sigma$ :
\[
z^{(\ell)}_j = \sum_k w^{(\ell)}_{jk} a^{(\ell-1)}_k + b^{(\ell)}_j \qquad \text{et} \qquad a^{(\ell)}_j = \sigma\big(z^{(\ell)}_j\big)
\]
avec $a^{(0)} = x = (x_1, x_2, \dots, x_d)^T \in \mathbb{R}^d$ en entrée et $a^{(L)} = \hat{y}$ en sortie. Pour ajuster ces paramètres en minimisant le coût $C$, la rétropropagation (\emph{Backpropagation}) fait remonter l'erreur locale $\delta^{(\ell)}_j = \frac{\partial C}{\partial z^{(\ell)}_j}$ de la sortie vers l'entrée selon les quatre équations du cours :
\begin{align}
\delta^{(L)}_j &= \frac{\partial C}{\partial a^{(L)}_j}\,\sigma'\big(z^{(L)}_j\big) && \text{(EQ1 : erreur sur la couche de sortie $L$)} \\
\delta^{(\ell)}_j &= \left(\sum_k w^{(\ell+1)}_{kj}\,\delta^{(\ell+1)}_k\right)\sigma'\big(z^{(\ell)}_j\big) && \text{(EQ2 : rétropropagation de $\ell+1$ vers $\ell$)} \\
\frac{\partial C}{\partial b^{(\ell)}_j} &= \delta^{(\ell)}_j \quad \text{et} \quad \frac{\partial C}{\partial w^{(\ell)}_{jk}} = a^{(\ell-1)}_k\,\delta^{(\ell)}_j && \text{(EQ3 et EQ4 : gradients des biais et des poids)}
\end{align}
Les poids et les biais sont ensuite mis à jour par \textbf{descente de gradient stochastique (SGD)} par mini-batchs, complétée par un terme d'inertie (\emph{momentum} $\mu = 0{,}9$).

\paragraph{Pourquoi un CNN en bonus et pourquoi Keras ?}
Dans le MLP, la première couche reçoit directement le vecteur aplati de $d=1\,024$ pixels. Impossible alors pour le réseau de savoir que deux pixels étaient voisins verticalement dans l'image d'origine. Pour contourner ce problème, des recherches hors-cours nous ont amenés à tester un réseau convolutif (CNN) en modèle bonus afin de conserver l'image en 2D ($32\times32$). Le principe reste intuitif : avant la couche dense finale, de petits filtres carrés ($3\times3$) balaient l'image 2D pour détecter localement les traits, coins et boucles, tout en réduisant peu à peu la taille de la grille (\emph{MaxPooling}).

En TP, seul un perceptron binaire simple avait été codé de zéro en \texttt{NumPy}, le reste reposant sur \texttt{scikit-learn}. Comme \texttt{scikit-learn} ne gère pas les couches convolutives et que les programmer à la main en \texttt{NumPy} aurait pris des semaines, ce bonus a été écrit avec \texttt{Keras}. L'architecture enchaîne deux couches de filtres $3\times3$ (32 puis 64 filtres), une couche dense de 128 neurones (\texttt{dropout} $= 0{,}2$) et 15 neurones de sortie (\emph{Softmax}), entraînés par SGD ($\eta = 0{,}02$, \texttt{batch\_size} $= 32$) selon les mêmes équations de rétropropagation.

%======================== RESULTATS ========================
\section{Résultats}\label{sec:res}

Le tableau~\ref{tab:res} regroupe les résultats obtenus sur les 3\,000 images de test.

\begin{table}[H]
\centering
\caption{Performances sur l'ensemble de test pour chaque modèle.}
\label{tab:res}
\begin{tabular}{lcccc}
\toprule
\textbf{Modèle} & \textbf{Images d'entraînement} & \textbf{Temps d'entraînement final} & \textbf{Accuracy test} & \textbf{$F_1$ macro} \\
\midrule
Random Forest (1D) & 12\,000 & 33,0 s & 81,20\,\% & 81,1\,\% \\
MLP (1D) & 12\,000 & 14,9 s & 87,67\,\% & 87,8\,\% \\
SVM (1D) & 12\,000 & 35,1 s & 87,97\,\% & 87,9\,\% \\
CNN (2D) & 12\,000 & 175,0 s & \textbf{97,60\,\%} & \textbf{97,6\,\%} \\
\bottomrule
\end{tabular}
\end{table}

Le \textbf{CNN} arrive nettement en tête avec \textbf{97,60\,\%} d'accuracy ($+9{,}63$ points devant le SVM). Parmi les trois modèles travaillant sur vecteurs 1D, le SVM (87,97\,\%) et le MLP (87,67\,\%) se tiennent à seulement 0,30 point d'écart, tandis que le Random Forest reste distancé à 81,20\,\%. Pour les quatre approches, l'accuracy et le $F_1$ macro sont quasiment identiques.

À l'échelle des classes individuelles, les modèles 1D montrent pourtant des écarts marqués entre précision et rappel. Le Random Forest illustre bien ce phénomène : sur la classe 9, la précision atteint 0,839 mais le rappel tombe à 0,650 (plus d'un tiers des 9 ne sont pas détectés) ; la classe 100 suit la même tendance (précision de 0,836 contre un rappel de 0,690). La classe 10 présente le défaut inverse avec beaucoup de faux positifs (précision à 0,716 pour un rappel à 0,795). Les matrices de confusion correspondantes sont présentées aux figures~\ref{fig:rf} à~\ref{fig:conf_cnn}.

\begin{figure}[H]
\centering
\includegraphics[width=0.55\textwidth]{confusion_random_forest.png}
\caption{Matrice de confusion du Random Forest sur le jeu de test (accuracy : 81,20\,\%). Les classes correspondent aux valeurs numériques.}
\label{fig:rf}
\end{figure}

\begin{figure}[H]
\centering
\includegraphics[width=0.55\textwidth]{confusion_mlp.png}
\caption{Matrice de confusion du MLP sur le jeu de test (accuracy : 87,67\,\%).}
\label{fig:mlp}
\end{figure}

\begin{figure}[H]
\centering
\includegraphics[width=0.55\textwidth]{confusion_svm.png}
\caption{Matrice de confusion du SVM sur le jeu de test (accuracy : 87,97\,\%). C'est le meilleur de nos trois modèles sur vecteurs 1D.}
\label{fig:svm}
\end{figure}

\begin{figure}[H]
\centering
\includegraphics[width=0.55\textwidth]{confusion_cnn.png}
\caption{Matrice de confusion du CNN 2D optimisé sur le jeu de test (accuracy : 97,60\,\%).}
\label{fig:conf_cnn}
\end{figure}

%======================== ANALYSE CRITIQUE ========================
\section{Analyse critique}\label{sec:crit}

\subsection{Comparaison des trois modèles de base (en 1D)}
Entre le SVM (87,97\,\%) et le MLP (87,67\,\%), la différence se joue à 9 images près sur les 3\,000 du jeu de test. La vraie cassure sépare ces deux approches du Random Forest (81,20\,\%).

\begin{enumerate}
\item \textbf{Plafonnement du Random Forest (81,20\,\%).} Chaque arbre découpe l'espace en posant des seuils sur des coordonnées isolées ($x_i \le \theta$). Aucune relation de voisinage 2D n'est prise en compte. Il suffit qu'un scripteur décale son chiffre de trois pixels sur le côté pour que d'autres colonnes du vecteur s'allument, mettant en défaut les règles apprises.

\item \textbf{SVM contre MLP.} Le noyau RBF compare les images globalement via la distance euclidienne $\exp(-\gamma\|x-x'\|^2)$ et s'appuie sur les vecteurs de support proches de la frontière de décision. Avec $C=10$, son score sur les 12\,000 images d'entraînement grimpe d'ailleurs à 99,18\,\%, signe d'un fort surapprentissage, même si la maximisation de la marge maintient un bon score en test (87,97\,\%). Le MLP, lui, combine tous les pixels dans ses couches cachées ($z^{(\ell)}_j$) mais comme il travaille sur le vecteur aplati, il reste sensible aux décalages ou à l'inclinaison du caractère, tout en s'entraînant plus de deux fois plus vite (14,9~s contre 35,1~s pour le SVM).
\end{enumerate}

\subsection{Analyse des confusions géométriques}
Le tableau~\ref{tab:conf} détaille les $F_1$-scores sur les dix classes les plus délicates à séparer.

\begin{table}[H]
\centering
\caption{$F_1$ par classe sur l'ensemble de test (200 images par classe), pour les classes les plus difficiles.}
\label{tab:conf}
\begin{tabular}{lcccc}
\toprule
\textbf{Classe} & \textbf{Random Forest} & \textbf{MLP} & \textbf{SVM} & \textbf{CNN (2D)} \\
\midrule
\zh{二} (2) & 0,810 & 0,913 & 0,846 & 0,972 \\
\zh{三} (3) & 0,821 & 0,915 & 0,862 & 0,980 \\
\zh{六} (6) & 0,767 & 0,846 & 0,854 & 0,985 \\
\zh{七} (7) & 0,797 & 0,878 & 0,868 & 0,985 \\
\zh{九} (9) & 0,732 & 0,766 & 0,836 & 0,977 \\
\zh{十} (10) & 0,754 & 0,832 & 0,857 & 0,945 \\
\zh{百} (100) & 0,756 & 0,831 & 0,835 & 0,952 \\
\zh{千} (1\,000) & 0,756 & 0,793 & 0,822 & 0,940 \\
\zh{万} (10\,000) & 0,790 & 0,826 & 0,839 & 0,973 \\
\zh{亿} ($10^8$) & 0,811 & 0,888 & 0,931 & 0,982 \\
\bottomrule
\end{tabular}
\end{table}

Sur les modèles 1D, l'essentiel des confusions touche les puissances de dix (10 à 10\,000) et le chiffre 9.

La confusion la plus fréquente oppose \textbf{\zh{十} (10) et \zh{千} (1\,000)}. Seule une courte barre oblique au sommet de \zh{千} le distingue de la croix de \zh{十} ; après passage en $32\times32$, ce détail se fond souvent dans la barre horizontale dès que le feutre est épais (les modèles 1D plafonnent entre 0,754 et 0,857 de $F_1$, contre 0,945 et 0,940 pour le CNN). Un second bloc d'erreurs concerne les caractères à traits horizontaux empilés comme \textbf{\zh{二} (2) et \zh{三} (3)}, où le MLP devance déjà le SVM (0,913 et 0,915 contre 0,846 et 0,862) avant que les filtres $3\times3$ du CNN ne portent ces scores à 0,972 et 0,980. Enfin, \textbf{\zh{九} (9), \zh{七} (7) et \zh{亿} ($10^8$)} partagent tous les trois un crochet final courbé vers la droite très variable selon l'écriture : pire classe du Random Forest (0,732) et du MLP (0,766), le chiffre 9 remonte à 0,977 avec le CNN.

\subsection{Apport du modèle bonus (CNN)}\label{sec:apport_cnn}
Passer des vecteurs 1D au CNN fait grimper l'accuracy de 87,97\,\% (SVM) à \textbf{97,60\,\%}. Le nombre d'erreurs sur les 3\,000 images de test est divisé par cinq : 72 erreurs contre 361 pour le SVM, en 175,0 secondes d'entraînement sur CPU. Sur le tableau~\ref{tab:conf}, le CNN arrive en tête sur \textbf{chacune des 15 classes}.

La raison tient directement au maintien de la grille 2D ($32\times32$). Pour le MLP comme pour le SVM, deux pixels voisins verticalement se retrouvent séparés de 32 cases dans le vecteur $\mathbb{R}^{1024}$ et sont traités comme deux variables sans lien. Avec ses filtres $3\times3$ qui balaient l'image, le CNN repère un même trait ou croisement même si le scripteur l'a tracé un peu plus haut ou plus à droite. Sur les caractères chargés, le saut est net : entre le SVM et le CNN, le $F_1$ passe de 0,857 à 0,945 sur \zh{十} (10), de 0,835 à 0,952 sur \zh{百} (100), et de 0,839 à 0,973 sur \zh{万} (10\,000).

\subsection{Limites et pistes d'amélioration}
Le temps de calcul sur CPU nous a imposé plusieurs compromis. Tester les 12 configurations du CNN sur les 12\,000 images aurait pris des heures ; il a fallu limiter cette recherche à 3\,000 images et 12 époques (194,0~s), et évaluer l'ensemble sur un seul découpage (\texttt{seed = 42}). La réduction en $32\times32$ pixels — choisie au départ pour alléger les modèles 1D — gomme aussi les traits fins de \zh{千} (1\,000) et \zh{十} (10), qui restent nos deux seules classes sous 0,95 de $F_1$ (0,940 et 0,945). Relancer le CNN directement sur les images $64\times64$ d'origine avec de légères rotations aléatoires permettrait sans doute de corriger ces dernières erreurs.

%======================== CONCLUSION ========================
\section{Conclusion}\label{sec:concl}

Sur les images Chinese MNIST aplaties en vecteurs de 1\,024 pixels, le SVM à noyau RBF (87,97\,\%) et le MLP (87,67\,\%) arrivent presque à égalité sous la barre des 88\,\%, loin devant le Random Forest (81,20\,\%) qui traite chaque pixel isolément. Notre exploration d'un réseau convolutif 2D (CNN) en modèle bonus, entraîné par descente de gradient stochastique avec \emph{momentum}, porte l'accuracy de test à \textbf{97,60\,\%} ($+9{,}63$ points sur le SVM).

Le travail direct sur les matrices 2D élimine l'essentiel des confusions géométriques : les 15 classes dépassent désormais 0,940 de $F_1$-score. Il ne reste que 72 erreurs sur les 3\,000 images de test (2,40\,\%), principalement entre \zh{千} (1\,000) et \zh{十} (10) dont les traits se touchent après réduction en $32\times32$ pixels. Conserver la résolution d'origine de $64\times64$ pixels serait la première piste à tester pour aller plus loin.

\end{document}
